\hsection{Chapter 0}
\sol{1} $60$: $1,2,3,4,5,6,10,12,15,20,30,60$. $61$: $1,61$ (prime). $62$: $1,2,31,62$. $63$: $1,3,7,9,21,63$. $64$: $1,2,4,8,16,32,64$. $65$: $1,5,13,65$. Only $64=8^2$ is a square (a square has an odd number of divisors; $64$ has seven).

\sol{3} $\frac16-\frac17=\frac{7-6}{42}=\frac1{42}$. In general $\frac1p-\frac1q=\frac{q-p}{pq}$, which equals $\frac1{pq}$ exactly when $q-p=1$: $q=p+1$ (and $p,q\ne0$), e.g.\ $\frac12-\frac13=\frac16$.

\sol{6} $x\mapsto x^2$: a parabola through $(-1,1),(0,0),(1,1),(2,4)$, symmetric in the $y$-axis, minimum $0$ at $x=0$. $x\mapsto\sqrt x$ on $[0,\infty)$: starts at $(0,0)$ with a vertical tangent, through $(1,1),(4,2)$, increasing and bending downwards (it is the mirror image of the right half of the parabola in the line $y=x$).

\sol{7} $H(x)=0$ for $x<0$ and $H(x)=1$ for $x\ge0$ (the closed dot in the dictaat's graph is at $(0,1)$; if it were on the lower branch, use $x\le0$ and $x>0$).

\sol{8} $y=3-\frac1{2+x}\Rightarrow\frac1{2+x}=3-y\Rightarrow2+x=\frac1{3-y}$ (for $y\ne3$) $\Rightarrow x=\frac1{3-y}-2$. Swapping letters: $f_{\mathrm{inv}}(x)=\frac1{3-x}-2$; $f$ is bijective from $\R\setminus\{-2\}$ onto $\R\setminus\{3\}$.

\sol{12} $1-\cos^2x=\sin^2x$, so $\sqrt{1-\cos^2x}=\sqrt{\sin^2x}=|\sin x|$. Sanity check $x=-\frac\pi6$: $\cos^2=\frac34$, $\sqrt{1-\frac34}=\frac12$, and $|\sin(-\frac\pi6)|=|-\frac12|=\frac12$ \checkmark\ (the answer $\sin x=-\frac12$ would be wrong).

\sol{13} $(p\circ p)(x)=1+(1+x)=2+x$; $(p\circ q)(x)=1+x^2$; $(q\circ p)(x)=(1+x)^2$; $(q\circ q)(x)=x^4$.

\sol{14} Many: $f(x)=x$; $f(x)=a-x$ for any constant $a$; $f(x)=\frac1x$ and $f(x)=\frac ax$ ($x\ne0$); $f(x)=\frac{x+1}{x-1}$ (check: $f(f(x))=\frac{\frac{x+1}{x-1}+1}{\frac{x+1}{x-1}-1}=\frac{2x}{2}=x$), more generally $f(x)=\frac{ax+b}{cx-a}$. Such functions are their own inverse: their graph is symmetric in the line $y=x$.

\sol{16} (a) Even: $x^2$, $x^4$, $|x|$, $\cos x$, constants. Odd: $x$, $x^3$, $\frac1x$, $\sin x$. (b) $x+1$, $e^x$, $x^2+x$. (c) If $f$ is both then $f(x)=f(-x)=-f(x)$, so $f(x)=0$: only the zero function (on any domain symmetric about $0$).

\sol{18} If $0$ is in the domain of an odd $f$, then $f(0)=-f(-0)=-f(0)$, so $f(0)=0$: the claim is right. It fails only when $0$ is not in the domain: $f(x)=\frac1x$ is odd and has no value at $0$. For even functions nothing of the kind holds: $f(0)$ can be anything ($\cos0=1$). (What is true: a differentiable even function has $f'(0)=0$, since $f'$ is odd.)

\sol{20(b)} Order the divisor: $2x^3+x$. $6x^4\div2x^3=3x$; subtract $3x(2x^3+x)=6x^4+3x^2$: remainder $-4x^3-4x^2+6x-12$. $-4x^3\div2x^3=-2$; subtract $-2(2x^3+x)=-4x^3-2x$: remainder $-4x^2+8x-12$, of degree $2<3$. So $\dfrac{6x^4-4x^3-x^2+6x-12}{2x^3+x}=3x-2+\dfrac{-4x^2+8x-12}{2x^3+x}$. Check: $(3x-2)(2x^3+x)+(-4x^2+8x-12)=6x^4+3x^2-4x^3-2x-4x^2+8x-12=6x^4-4x^3-x^2+6x-12$ \checkmark.

\sol{21(b)} The common factor is $5x^2(x-14)$: $5x^2(x-14)(9x^2-10x+6)$. The quadratic has discriminant $100-216<0$, so it has no real zeros and cannot be factorised further over $\R$. Check: $5x^2(x-14)\cdot9x^2=45x^4(x-14)$, etc.

\sol{21(c)} $x-x^3+x^5-x^7=x(1-x^2+x^4-x^6)=x\big((1-x^2)+x^4(1-x^2)\big)=x(1-x^2)(1+x^4)=x(1-x)(1+x)(1+x^4)$. Check: $(1-x^2)(1+x^4)=1+x^4-x^2-x^6$ \checkmark. ($1+x^4$ has no real zeros; over $\R$ it still splits as $(x^2+\sqrt2x+1)(x^2-\sqrt2x+1)$, but that is beyond ``simple'' factorising.)

\sol{22} Square with diagonal $1$: the diagonal makes angles $\frac\pi4$ with the sides; the sides $s$ satisfy $s^2+s^2=1$, so $s=\frac1{\sqrt2}=\frac12\sqrt2$. Equilateral triangle with side $1$: all angles $\frac\pi3$; its altitude splits it into two right triangles with angles $\frac\pi6,\frac\pi3,\frac\pi2$, hypotenuse $1$, short side $\frac12$ and altitude $\sqrt{1-\frac14}=\frac12\sqrt3$. These give the table of Oefening 23.

\sol{24} (a) The point at angle $t$ is $(\cos t,\sin t)$: $\cos t$ is its horizontal coordinate (the segment on the $x$-axis), $\sin t$ its vertical coordinate. The vertical tangent line $x=1$ meets the extended radius at height $h$; the triangle $(0,0),(1,0),(1,h)$ is similar to $(0,0),(\cos t,0),(\cos t,\sin t)$, so $h/1=\sin t/\cos t=\tan t$. (b) $\sin(\pi+t)=-\sin t$, $\sin(\pi-t)=+\sin t$, $\cos(\pi+t)=-\cos t$, $\cos(\pi-t)=-\cos t$, $\tan(\pi+t)=+\tan t$, $\tan(\pi-t)=-\tan t$ (rotating by $\pi$ negates both coordinates; reflecting in the $y$-axis negates only the first). (c) Reflection in the $x$-axis ($t\mapsto-t$) negates the $y$-coordinate and keeps the $x$-coordinate: $\sin(-t)=-\sin t$, $\cos(-t)=\cos t$. Reflection in the line $y=x$ ($t\mapsto\frac\pi2-t$) swaps the coordinates: $\sin(\frac\pi2-t)=\cos t$. Pythagoras in the triangle with legs $\cos t,\sin t$ and hypotenuse $1$: $\sin^2t+\cos^2t=1$.

\sol{27} (a) The velocity of a point moving on a circle is tangent to the circle, and the tangent is perpendicular to the radius $OP$ (a chord $PQ$ with $Q\to P$ becomes perpendicular to $OP$ by symmetry of the isosceles triangle $OPQ$). (b) A vector of length $1$ at $P$, perpendicular to $OP$, pointing counterclockwise. (c) Rotating $OP=(\cos t,\sin t)$ over $+\frac\pi2$ gives $(-\sin t,\cos t)$: the endpoint is $(-\sin t,\cos t)$. (d) Velocity is the derivative of position: $\frac d{dt}(\cos t,\sin t)=(-\sin t,\cos t)$, so $\cos't=-\sin t$ and $\sin't=\cos t$.

\sol{29} (a) ${}^h\!\log x=\dfrac{{}^g\!\log x}{{}^g\!\log h}$ for all $x>0$, so ${}^g\!\log x={}^g\!\log h\cdot{}^h\!\log x$: the two functions differ by the constant factor ${}^g\!\log h$. (b) The factor is ${}^g\!\log h$ (equivalently $1/{}^h\!\log g$). (c) Taking $x=g$: ${}^h\!\log g=\dfrac{{}^g\!\log g}{{}^g\!\log h}=\dfrac1{{}^g\!\log h}$, so ${}^g\!\log h\cdot{}^h\!\log g=1$. True statements in the list: ``you can always compute it if you know $g$ and $h$'', ``it exists provided $g$ and $h$ are positive (and $\ne1$, so that they can be bases)'', ``it is $1$''. It does not depend on which base the notation $\log$ refers to, and it is not $0$.
